\section{円の媒介変数表示}\label{節:円の有理パラメータ}

前節の2式は
\[
X=\frac{1-t^2}{1+t^2},\qquad
Y=\frac{2t}{1+t^2}
\]
と書くと,\,$X^2+Y^2=1$を満たす.正弦と余弦を使ったのだから当然ともいえるが,この式は三角関数なしでも導ける.

図\ref{図:円の有理パラメータ}のように,単位円上の点$A=(-1,0)$を固定し,傾き$t$の直線$Y=t(X+1)$を引く.円との交点を求めるため代入すると
\[
\begin{aligned}
X^2+t^2(X+1)^2&=1,\\
(X+1)\{(1+t^2)X+(t^2-1)\}&=0.
\end{aligned}
\]
1つの解$X=-1$は,固定した点$A$である.もう1つの交点$P$は,上の$X,Y$で表される.\,$t$を動かして点を表すので,\,$t$を\textbf{媒介変数},この表示を円の\textbf{有理パラメータ表示}という.

\begin{figure}[htbp]
\centering
\begin{tikzpicture}[x=2.25cm,y=2.25cm,>=Stealth]
 \draw[->] (-1.4,0)--(1.45,0) node[right]{$X$};
 \draw[->] (0,-1.2)--(0,1.35) node[above]{$Y$};
 \draw[thick,tealblack] (0,0) circle(1);
 \draw[thick,blueblack] (-1.23,-0.115)--(1.15,1.075);
 \draw[dashed] (0.6,0)--(0.6,0.8)--(0,0.8);
 \draw[blueblack] (0,0)--(0.6,0.8);
 \draw[->] (0.26,0) arc[start angle=0,end angle=53.1301,radius=0.26];
 \node at (0.35,0.17) {$x$};
 \draw[->] (-0.65,0) arc[start angle=0,end angle=26.5651,radius=0.35];
 \node at (-0.42,0.15) {$\frac x2$};
 \fill (-1,0) circle(1.2pt) node[below left]{$A=(-1,0)$};
 \fill (0.6,0.8) circle(1.2pt);
 \draw[thin,gray] (0.63,0.79)--(1.04,0.56);
 \node[anchor=west] at (1.04,0.52) {$P=(\cos x,\sin x)$};
 \node[above right] at (0.97,0.985) {$Y=t(X+1)$};
 \node[below left] at (0,0) {$O$};
\end{tikzpicture}
\caption{円上の固定点から引いた直線の傾きで,もう1つの交点を表す.図は$t=\frac12$の場合.}
\label{図:円の有理パラメータ}
\end{figure}

$P=(X,Y)\ne A$からは$t=\frac{Y}{X+1}$と戻せるので,有限な実数$t$で,円上の$A$以外のすべての点を一度ずつ表せる.\,$P=(\cos x,\sin x)$なら
\[
t=\frac{\sin x}{1+\cos x}=\tan\frac x2
\]
である.最後の等式は2倍角の公式から確かめられる.図の上半円では,三角形$AOP$の角を見ても,直線$AP$が正の$X$方向となす角は$\frac x2$になる.実際,\,$OA=OP=1$より三角形$AOP$は二等辺三角形であり,頂角が$\angle AOP=\pi-x$なので,底角$\angle OAP$は$\frac{\pi-(\pi-x)}{2}=\frac x2$である.\,$A$から$O$へ向かう向きは正の$X$方向なので,この底角が求める角である.半角が現れる理由を,円の上の角としても読めた.表せなかった点$A$は,ちょうど$x=(2k+1)\pi$に対応する.

\subsection{ピタゴラス数}

$t=\frac nm$を整数$m,n$で表すと,
\[
X=\frac{m^2-n^2}{m^2+n^2},\qquad
Y=\frac{2mn}{m^2+n^2}.
\]
従って
\begin{equation}\label{式:ピタゴラス数}
(m^2-n^2)^2+(2mn)^2=(m^2+n^2)^2
\end{equation}
である.\,$m>n>0$なら三つの整数は正であり,直角三角形の三辺になる.\,$a^2+b^2=c^2$を満たす正の整数の組を\textbf{ピタゴラス数}という.

逆に整数の直角三角形$(a,b,c)$があれば,点$\left(\frac ac,\frac bc\right)$は単位円上にある.その傾きは
\[
t=\cfrac{\cfrac bc}{1+\cfrac ac}=\frac{b}{c+a}
\]
と有理数になる.これを$\frac nm$と表せば,上の公式が元の三辺と同じ比を与える.従って,この表示は偶然いくつかの例を作るだけでなく,すべての整数の直角三角形の辺の比を表す.得られた三整数に共通因数があれば,それで割ればよい.

三辺の最大公約数が$1$のピタゴラス数を,\textbf{原始ピタゴラス数}という.整数$m>n>0$が互いに素で,一方が偶数,他方が奇数ならば,式\eqref{式:ピタゴラス数}は原始ピタゴラス数を与える.逆に,すべての原始ピタゴラス数は,直角を挟む二辺の交換を除いてこの形で表せる.この主張を,整数の素因数分解の一意性を使って確かめよう.

まず上の条件で作った三整数に,共通の素因数$p$があると仮定する.\,$m^2-n^2$と$m^2+n^2$はともに奇数なので,\,$p$は奇素数である.この二数の和と差から,\,$p$は$2m^2$と$2n^2$を割る.従って$p$は$m,n$の両方を割り,互いに素という条件に反する.共通の素因数がないので,三整数の最大公約数は$1$である.

逆に原始ピタゴラス数$(a,b,c)$を考える.二辺$a,b$がともに偶数なら原始ではなく,ともに奇数なら$a^2+b^2$を$4$で割った余りが$2$となり,平方数$c^2$にはなれない.従って一方だけが偶数であり,必要なら辺を交換して$b$を偶数,\,$a$を奇数とする.このとき$c$も奇数である.また$a,c$に共通の素因数があれば,\,$b^2=c^2-a^2$,従って$b$もその素因数で割れ,原始という条件に反する.よって$a,c$は互いに素である.

ここで
\[
A=\frac{c+a}{2},\qquad B=\frac{c-a}{2}
\]
と置くと,\,$A>B>0$の整数であり,
\[
AB=\frac{c^2-a^2}{4}=\left(\frac b2\right)^2.
\]
$A,B$の共通因数は$A+B=c$と$A-B=a$も割るので,\,$A,B$は互いに素である.互いに素な二数の積が平方数なら,素因数は片方にしか現れず,その指数はすべて偶数になる.従って$A=m^2$,\,$B=n^2$と書ける.和と差を戻せば
\[
a=m^2-n^2,\qquad b=2mn,\qquad c=m^2+n^2.
\]
$A,B$が互いに素だから$m,n$も互いに素であり,\,$c$が奇数だから$m,n$の偶奇は異なる.これで逆向きも示せた.円の表示が辺の比を与えることに加えて,約分の要らない整数を作る条件まで分かった.

\begin{example}{ピタゴラス数}
$t=\frac23$からピタゴラス数を作れ.また三辺が$8,15,17$の三角形に対応する傾きを求めよ.
\end{example}

$m=3,n=2$なので
\[
m^2-n^2=5,\qquad2mn=12,\qquad m^2+n^2=13.
\]
確かに$5^2+12^2=13^2$である.後半では$P=\left(\frac{15}{17},\frac8{17}\right)$と選ぶと
\[
t=\frac8{17+15}=\frac14.
\]
$m=4,n=1$を公式へ戻せば$(15,8,17)$となる.二つの直角を挟む辺のどちらを$X$座標にするかで,対応する傾きは変わる.



%
